%%%%%%%%%%%%%%%%%%%%%%%%%%%%%%%%%%%%%%%%%%%%%%%%%%%%%%%%%%%%
\section[Bilder, Grafiken und Diagramme]{Bilder, Grafiken und Diagramme}
\label{sec:Bilder}
%%%%%%%%%%%%%%%%%%%%%%%%%%%%%%%%%%%%%%%%%%%%%%%%%%%%%%%%%%%%
%
Bei Einbindung von Grafiken sind zwei Fälle zu unterscheiden:
\begin{itemize*}
  \item reguläre \index{Bild!Binär-}Bilder in einem Binärformat
	      (\texttt{PNG}, \texttt{TIFF}, \texttt{JPG}, \texttt{PDF}, \etc)
	\item \index{Bild!Vektor-}Grafiken, die im \gls{gls:tikz}-Quellcode vorliegen 
\end{itemize*}

Grundlegender Unterschied bei der Einbindung \enquote{regulärer} Bilder
und \gls{gls:tikz}-Bilder ist, dass Binärformatgrafiken mit
\lc{includegraphics\{...\}} eingebunden werden,
während \gls{gls:tikz}-Grafiken mit \lc{input\{...\}}
eingebunden und von \printkeyword{latex} mitkompiliert werden.


%%%%%%%%%%%%%%%%%%%%%%%%%%%%%%%%%%%%%%%%%%%%%%%%%%%%%%%%%%%%
\subsection[Floats]{\index{Float}\index{Bild!Float}Floats}%
\label{sec:Floats}
%%%%%%%%%%%%%%%%%%%%%%%%%%%%%%%%%%%%%%%%%%%%%%%%%%%%%%%%%%%%
%
Üblicherweise werden Bilder, Tabellen und Codelistings in Fließumgebungen (floats) gesetzt,
damit \LaTeX\ sie frei im Text herumbewegen und dadurch geschickt positionieren kann.
Bei Bildern heißt die entsprechende Float-Umgebung \printkeyword{figure}, 
bei Tabellen \printkeyword{table}, bei Codelistings \printkeyword{lstlisting}.
Die Positionierung kann durch Angabe von Buchstaben
\printkeyword{h}, \printkeyword{t}, \printkeyword{b} und \printkeyword{p}
beeinflusst werden
(\enquote{\textbf{h}ere}, \enquote{\textbf{t}op}, \enquote{\textbf{b}ottom}, \enquote{extra \textbf{p}age}).
Eine genaue Auswirkung der Parameter
\printkeyword{h}, \printkeyword{t}, \printkeyword{b} und \printkeyword{p}
\bzw ihrer Kombinationen auf die Bildplatzierung ist jedoch nicht intuitiv.
Um diese zu verstehen, empfiehlt sich die Lektüre der Beschreibung
von Frank Mittelbach
\href{https://tex.stackexchange.com/questions/39017/how-to-influence-the-position-of-float-environments-like-figure-and-table-in-lat/39020#39020}{auf stackexchange.com}.%
\footnote{\url{https://tex.stackexchange.com/questions/39017/how-to-influence-the-position-of-float-environments-like-figure-and-table-in-lat/39020#39020}}

So versucht \LaTeX{} bei angegebener \printkeyword{t}- oder \printkeyword{b}-Option,
das Bild \bzw die Tabelle am oberen oder unteren Rand der Seite zu platzieren,
wodurch es passieren kann, dass der Absatz durch die Abbildung \etc
mitten in einem Satz aufgetrennt wird, was im professionellen Buchdruck unerwünscht ist.
Auch bei Einbindung von Floats
%mitten in einem Textabsatz
mit Hilfe der \printkeyword{h}-Option
platziert \LaTeX{} diese nicht unbedingt direkt an der Stelle der Einbindung im Quellcode.
%sondern zwischen zwei Zeilen in der Nähe der Einbindungsstelle, wobei auf Einhaltung
%möglichst gleichmäßiger Wortabstände im Absatz geachtet wird.
%Somit kann es dabei ebenfalls zu einer unerwünschten Auftrennung des Textes
%mitten in einem Satz kommen.
%
%Auch bei der Platzierung von Bildern und Tabellen mit Hilfe der \printkeyword{h}-Option
%in einem eigenen Absatz zwischen zwei Textabsätzen kann es zum unerwünschten Verhalten kommen.
%Wenn der verbleibende Platz die Platzierung eines Floats an der erwünschten Stelle nicht erlaubt,
%kann es schnell passieren, dass das betroffene Float auf die nächste Seite verschoben werden muss.
%Dabei wandelt \LaTeX{} die \printkeyword{h}-Option oft eigenwillig in die \printkeyword{ht}-Option um
%und verschiebt das Bild an den oberen Rand der Seite,
%%(beim Kompilieren des Quellcodes erhält man hierzu Feedback in Form einer Warnung), 
%wobei die oben beschriebene unerwünschte Auftrennung des Textes mitten in einem Satz entstehen kann.
%Muss ein Float soweit verschoben werden, dass das nächste Float von \LaTeX{} mitangefasst werden müsste, 
%wird das betroffene Float und alle Folgefloats unerwünschterweise ans Ende des Kapitels verschoben
%(genauer gesagt, an die Stelle, wo die nächste \lstinline|\clearpage|-Anweisung kommt).

Eine Abhilfe schafft eine konsequente Einbindung von Bildern und Tabellen mit Hilfe der \printkeyword{H}-Option,
%deren Nutzung in dieser Vorlage durch die Einbindung des \pkg{float}-Pakets aktiviert wurde und
welche eine Platzierung von Floats an der spezifizierten Stelle erzwingt.
Die Einbindung sollte als ein eigener Absatz zwischen zwei Textabsätzen erfolgen.

\myexcl{Wichtig!}
Wenn der verbleibende Platz die Platzierung eines Floats an der erwünschten Stelle nicht erlaubt,
wird dieses auf die nächste Seite verschoben. Dasselbe gilt für den nachfolgenden Inhalt.
Als Nebenwirkung können dabei auf der vorhergehenden Seite riesige leere Flächen entstehen.
Falls solche Nebeneffekte auftreten, müssen die betroffenen Bilder- und Tabellendefinitionen
manuell im Quellcode solange zwischen den Absätzen herumgeschoben werden,
bis eine optisch annehmbare Platzierung erreicht wird.
Dafür empfiehlt es sich, die einzelnen Bild- und Tabellendefinitionen in Extra-Dateien auszulagern,
so dass nur noch jeweils eine einzige Zeile mit der \lc{input\{...\}}-Anweisung herumgeschoben werden muss.
Prinzipiell empfiehlt sich eine endgültige Platzierung der Bilder und Tabellen erst ganz am Schluss,
nachdem alle anderen Korrekturen durchgeführt worden sind.

Wichtig anzumerken ist, dass alle Zeilen innerhalb der \printkeyword{figure}"=
\bzw der \printkeyword{table}"=Umgebung
mit einem Prozentzeichen abzuschließen sind.
Ansonsten werden überflüssige Leerzeichen eingefügt,
was zu unerwünschten Nebenwirkungen führen kann.

%%%%%%%%%%%%%%%%%%%%%%%%%%%%%%%%%%%%%%%%%%%%%%%%%%%%%%%%%%%%
\subsection[Binärbilder]{\index{Bild!Binär-}Binärbilder}%
\label{sec:Binaerbilder}
%%%%%%%%%%%%%%%%%%%%%%%%%%%%%%%%%%%%%%%%%%%%%%%%%%%%%%%%%%%%
%
Ein Beispiel für die Einbindung eines Bildes im Binärformat zeigt \cref{lst:binary-image}.
%Ein Beispiel für die Einbindung eines Bildes im Binärformat ist in \cref{lst:binary-image} angeführt.
%Man beachte die abschließenden Prozent-Zeichen am Ende jeder Zeile!

\begin{latex}[caption={Einbindung einer Binärgrafik in LaTeX},label={lst:binary-image}]
\begin{figure}[H]%
  \centering%
  \includegraphics[width=\linewidth]{Bildpfad/Dateiname}%
  \caption[Kurzversion für das Abbildungsverzeichnis]{%
           Eine tolle sehr lange Abbildungsunterschrift}%
  \label{fig:my-binary-image}%
\end{figure}
\end{latex}

Die Angabe des Pfades kann sowohl absolut als auch relativ
zum Verzeichnis der Hauptdatei oder zu einem der Pfade angegeben werden,
die in der Datei \printfilepathquoted{./preambel/AllePfade.tex} definiert sind.
Diese Pfade werden in angegebener Reihenfolge durchsucht.
Dasselbe gilt für die Dateierweiterung.
Ist keine Erweiterung definiert und liegen mehrere Bilder mit gleichem Namen jedoch unterschiedlicher Dateierweiterung vor,
wird die Reihenfolge, die in der Datei \printfilepathquoted{AllePfade.tex} definiert ist, verwendet.

Der \gls{ac:KSP} empfiehlt, Bilder und Grafiken auf die Satzspiegelbreite zu skalieren.
Dies wird durch die Option \printkeyword{width=\bs linewidth} umgesetzt.

Das Ergebnis sieht man in \cref{fig:my-binary-image}.

\begin{figure}[H]%
  \centering%
  \includegraphics[width=\linewidth]{./images/examples/img-example-01.jpg}%
  \caption[Einbindung eines Bildes]{Beispiel für die Einbindung eines Binärbildes}%
  \label{fig:my-binary-image}%
\end{figure}

Mit Hilfe von Paket \pkg{subfig} \cite{Cochran2005} können Bilder auch in
\index{Abbildung|see{Bild}}\index{Bild!Unterabbildung}Unterabbildungen gesetzt
und sowohl als ganzes (vgl. \cref{fig:subfloat-example}) als auch einzeln (vgl.
\cref{fig:subfloat-example-01,fig:subfloat-example-02,fig:subfloat-example-03,fig:subfloat-example-04})
referenziert werden.

\begin{figure}[H]%
	\centering%
	\subfloat[La Savoureuse, Lepuix, Frankreich][La Savoureuse, Lepuix, Frankreich \linebreak \textcopyright{} Thomas Bresson]{%
		\label{fig:subfloat-example-01}%
		\includegraphics[width=0.49\linewidth]{./images/examples/subfloat-example-01.jpg}%
	}%
	\hfill%
	\subfloat[Bangkok, Thailand][Bangkok, Thailand \linebreak \textcopyright{} Prachanart Viriyaraks]{%
		\label{fig:subfloat-example-02}%
		\includegraphics[width=0.49\linewidth]{./images/examples/subfloat-example-02.jpg}%
	}%
	\\\vspace{3mm}%
	\subfloat[Wahkeena Falls, Lincoln Park, USA][Wahkeena Falls, Lincoln Park, USA \linebreak \textcopyright{} srslyguys]{%
		\label{fig:subfloat-example-03}%
		\includegraphics[width=0.49\linewidth]{./images/examples/subfloat-example-03.jpg}%
	}%
	\hfill%
	\subfloat[Nacionalni park Plitvička jezer, Kroatien][Nacionalni park Plitvička jezer, Kroatien \linebreak \textcopyright{} Roman Bonnefoy]{%
		\label{fig:subfloat-example-04}%
		\includegraphics[width=0.49\linewidth]{./images/examples/subfloat-example-04.jpg}%
	}%
	\caption[Bild mit Unterabbildungen]{Wasserfälle der Welt als Beispiel für Unterabbildungen}%
	\label{fig:subfloat-example}%
\end{figure}

Der Beispielcode dafür ist in \cref{lst:subfigures} dargestellt.

\begin{latex}[caption={Unterabbildungen in LaTeX},label={lst:subfigures}]
\begin{figure}[H]%
	\centering%
	\subfloat[Kurzversion][Unterbezeichnung 1]{%
		\label{fig:UnterAbb1}%
		\includegraphics[width=0.49\linewidth]{Bildpfad/Bild1}%
	}%
	\hfill%
	\subfloat[Kurzversion][Unterbezeichnung 2]{%
		\label{fig:UnterAbb2}%
		\includegraphics[width=0.49\linewidth]{Bildpfad/Bild2}%
	}%
	\\\vspace{3mm}%
	\subfloat[Kurzversion][Unterbezeichnung 3]{%
		\label{fig:UnterAbb3}%
		\includegraphics[width=0.49\linewidth]{Bildpfad/Bild3}%
	}%
	\hfill%
	\subfloat[Kurzversion][Unterbezeichnung 4]{%
		\label{UnterAbb4}%
		\includegraphics[width=0.49\linewidth]{Bildpfad/Bild4}%
	}%
\caption[Kurzversion]{Langversion der Bildunterschrift}%
\label{fig:MeinGanzesBild}%
\end{figure}
\end{latex}


%%%%%%%%%%%%%%%%%%%%%%%%%%%%%%%%%%%%%%%%%%%%%%%%%%%%%%%%%%%%
\subsection[TikZ-Grafiken]{\index{TikZ}\index{Bild!TikZ}\gls{gls:tikz}-Grafiken}%
\label{sec:TikZ}
%%%%%%%%%%%%%%%%%%%%%%%%%%%%%%%%%%%%%%%%%%%%%%%%%%%%%%%%%%%%
%
\Gls{gls:tikz} eignet sich hervorragend, um wissenschaftliche Zeichnungen,
Vektorgrafiken und \index{Diagramm}Diagramme direkt mithilfe von LaTeX
zu setzen, sodass die Schrift direkt zum restlichen Dokument passt.
Zu dem \pkg{tikz}-\gls{gls:pkg} und dem darauf aufsetzenden \pkg{PGFplots}-\gls{gls:pkg} gibt
es hervorragende Dokumentationen \parencites{Tantau2013}{Feuersaenger2014}.
Mit \gls{gls:tikz} und \gls{gls:pgfplots} lassen sich viele gute Sachen machen.
%
%\setcounter{figure}{10}

Der Code für die Einbindung einer \gls{gls:tikz}-Grafik steht in \cref{lst:tikz-figure}.
\begin{latex}[caption={Einbindung einer TikZ-Zeichnung in LaTeX},label={lst:tikz-figure}]
\begin{figure}[H]%
  \centering%
  \tikzsetnextfilename{TikZ-Bild}%
  \resizebox{\textwidth}{!}{%   <--- optionale Skalierung
    \input{./figures-src/TikZ-Bild.tex}%
  }%                            <--- optionale Skalierung
  \caption[Kurzversion für das Abbildungsverzeichnis]{%
           Eine tolle sehr lange Abbildungsunterschrift}%
  \label{fig:my-tikz-figure}%
\end{figure}
\end{latex}

Eine Skalierung auf die volle Seitenbreite oder ein vielfaches davon im Falle von Unterabbildungen kann bei Bedarf mit Hilfe der Anweisung
\lstinline|\resizebox{\textwidth}{!}{...}|
durchgeführt werden.

Das Kommando \lstinline|\tikzsetnextfilename{...}| ist nicht unbedingt notwendig,
 aber sehr zu empfehlen, da dies als Name für das temporäre Kompilat im Ordner
\printfilepathquoted{./figures-compiled/} genommen wird.
Dieser sollte gleich dem Namen der Quelldatei (ohne Endung) gewählt werden.
Ansonsten nimmt \texttt{pdflatex} eine hochlaufende Nummer als Dateiname,
was die Fehlersuche sehr erschwert.

Nachfolgend finden sich einige Beispiele für TikZ-Zeichnungen, nämlich
ein kommutatives Diagramm (\cref{fig:kpca}),
eine Übersicht über die KIT-Corporate-Identity-Farben (\cref{fig:kit-colors}),
ein Netzwerkkommunikationsgraph (\cref{fig:net-comm}),
einfache \index{Diagramm!Punkt-}Punktdiagramme (\cref{fig:ica})
und etwas aufwendigere Diagramme mit mehreren
\index{Achsensystem|see{Diagramm}}Achsensystemen (\cref{fig:pca}).

Vorzugsweise sollten für die Grafiken die KIT-Corporate-Identity-Farben verwendet werden,
die sowohl eine RGB"=Definition für die Darstellung online,
als auch eine CMYK"=Definition für den Offset-Druck haben.%
\footnote{Bei Erstellung der Manuskriptversionen
für die Online-Veröffentlichung \bzw für den Offset-Druck
ist auf die korrekte Einstellung der Option \printkeyword{useCMYKcolors}
in der Datei \printfilepathquoted{preambel/AlleSchalter.tex} zu achten
(\printkeyword{false} für die Online-Veröffentlichung und \printkeyword{true} für den Druck!}

\begin{figure}[H]%
	\centering%
  \tikzsetnextfilename{kpca}%
	%\resizebox{\textwidth}{!}{%
		\input{./figures-src/kpca.tex}%
	%}%
	\caption{Kommutative Diagramm mit TikZ}%
  \label{fig:kpca}%
\end{figure}

\begin{figure}[H]%
	\centering%
  \tikzsetnextfilename{kit-colors}%
	%\resizebox{\textwidth}{!}{%
		\input{./figures-src/kit-colors.tex}%
	%}%
	\caption{KIT-Corporate-Identity-Farben}%
  \label{fig:kit-colors}%
\end{figure}

\begin{figure}[H]%
	\centering%
  \tikzsetnextfilename{net-comm}%
	\resizebox{0.75\textwidth}{!}{%
		\input{./figures-src/net-comm.tex}%
	}%
	\caption{Netzwerkkommunikationsgraph mit TikZ}%
  \label{fig:net-comm}%
\end{figure}

\begin{figure}[H]%
	\centering%
	\subfloat[Ursprünglicher Merkmalsraum]{%
		\label{fig:ica-1}%
		\tikzsetnextfilename{ica-1}%
		\resizebox{0.45\textwidth}{!}{%
			\input{./figures-src/ica-1.tex}
		}%
	}%
	\hfill%
	\subfloat[Transienter Merkmalsraum (Nach Whitening, z.\,B. durch \gls{ac:PCA} inkl. Normalisierung)]{%
		\label{fig:ica-2}%
		\tikzsetnextfilename{ica-2}%
		\resizebox{0.45\textwidth}{!}{%
			\input{./figures-src/ica-2.tex}%
		}%
	}%
	\hfill%
	\subfloat[Transformierter Merkmalsraum]{%
		\label{fig:ica-3}%
		\tikzsetnextfilename{ica-3}%
		\resizebox{0.45\textwidth}{!}{%
			\input{./figures-src/ica-3.tex}%
		}%
}%
\caption{Diagramme mit TikZ direkt in LaTeX (hier: Die Schritte der \enquote{Independent component analysis})}%
\label{fig:ica}%
\end{figure}

\begin{figure}[H]
	\centering%
	\subfloat[Ungünstige Projektion]{%
		\label{fig:pca-1}%
		\tikzsetnextfilename{pca-1}%
		\resizebox{0.48\textwidth}{!}{%
			\input{./figures-src/pca-1.tex}%
		}%
	}\hfill%
	\subfloat[Zielführende Projektion]{%
		\label{fig:pca-2}%
		\tikzsetnextfilename{pca-2}%
		\resizebox{0.48\textwidth}{!}{%
			\input{./figures-src/pca-2.tex}%
		}%
	}%
	\caption{Aufwändiges Diagramm mit TikZ (hier: Probleme der \enquote{Principal component analysis})}%
	\label{fig:pca}%
\end{figure}